\documentclass[11pt]{amsart}
\IfFileExists{lmodern.sty}{\usepackage[T1]{fontenc}\usepackage{lmodern}}{}
\usepackage{amsmath,amssymb,amsthm}
\usepackage[hidelinks]{hyperref}
\setlength{\emergencystretch}{1.5em}
\IfFileExists{geometry.sty}{\usepackage[letterpaper,left=0.75in,right=0.75in,top=0.7in,bottom=0.8in]{geometry}}{\setlength{\textwidth}{7in}\setlength{\textheight}{9.4in}\setlength{\oddsidemargin}{-0.25in}\setlength{\evensidemargin}{-0.25in}\setlength{\topmargin}{-0.5in}}
\newtheorem{theorem}{Theorem}
\newtheorem{lemma}[theorem]{Lemma}
\newtheorem{corollary}[theorem]{Corollary}
\newtheorem{proposition}[theorem]{Proposition}
\newtheorem{remark}[theorem]{Remark}
\title[An $83/84$ approximation for a cubic three-occurrence Horn class]{An $83/84$ approximation for a cubic three-occurrence Horn class}
\author{Frank Vega}
\address{Information Physics Institute, 840 W 67th St, Hialeah, FL 33012, USA}
\email{vega.frank@gmail.com}
\date{November 22, 2026}
\subjclass[2020]{68Q25, 68W25, 90C27}
\keywords{approximate-computing, max-sat, independent-set-problem, lean4, np-hard, satisfiability, unique-games-conjecture}
\begin{document}
\begin{abstract}
We study an unweighted MAX-HORN-2SAT class, which is also a class of MAX 2-SAT instances, in which every Boolean variable has exactly three literal occurrences. For a simple cubic graph $H=(V,E)$ the formula $T_H$ satisfies
\[
  \operatorname{OPT}(T_H)=4|V|+|E|+\alpha(H),
\]
and every assignment of value $A$ is repaired in polynomial time to an independent set of size at least $A-(4|V|+|E|)$. This gives a polynomial-time $83/84$ approximation from a cubic independent-set approximation and a near-bipartite cut algorithm, an $L$-reduction from cubic MIS, and APX-completeness. For simple $d$-regular graphs the optimum is $(d+1)|V|+|E|+\alpha(H)$ and the factors are: $(3d+2)/(3d+3)$ from the all-true assignment, which is at least $83/84$ for $d\ge27$; $83/84$ from the Berman--Fujito guarantee for $d\ge22$ and, with a semidefinite cut rounding, for $d=21$; and $(21d+20)/(21d+21)$ from any $6/7$-approximate independent set. For arbitrary simple graphs the optimum is $n+3m+\alpha(H)$.

MAX-CUT on a $d$-regular graph $G$ reduces to MAX 2-SAT with the clauses $(x_u\vee x_v)$ and $(\neg x_u\vee\neg x_v)$ per edge, of value $m+\operatorname{cut}(S)$. The conflict graph $H(G)$ is $(d+1)$-regular and $\alpha(H(G))=m+\operatorname{maxcut}(G)$. A Goemans--Williamson cut gives on $T_{H(G)}$ the factor $1-\frac{1-\alpha_\varepsilon}{6d+12}\ge83/84$ for every $d\ge1$, where $\alpha_\varepsilon=\alpha_{\mathrm{GW}}-10^{-5}$. A cut guarantee $c_d$ on graphs of maximum degree $d$ gives $1-\frac{1-c_d}{6d+12}$, which exceeds $0.9977$ in expectation for $d=3$. A ratio on $T_{H(G)}$ and a cut guarantee are equivalent, so the cut recovered is the Goemans--Williamson or bounded-degree cut used, and the MAX-CUT guarantee is not improved.

The main consequence is that on the class $\{T_{H(G)}\}$ of three-occurrence Horn formulas obtained from MAX-CUT instances the proved factor $83/84$ exceeds the Unique Games thresholds $\approx0.94616$ and $\approx0.94017$ that hold for all instances of weighted MAX-HORN-2SAT and of MAX 2-SAT.
\end{abstract}
\maketitle
\section{Introduction}
MAX-HORN-2SAT maximizes the number of satisfied unary and binary Horn clauses. A Horn clause has at most one positive literal. The formula $T_H$ below is written with the opposite polarity: each of its clauses has at most one negative literal. Renaming every variable to its negation turns every clause of $T_H$ into a Horn clause, and an assignment satisfies a clause exactly when the complemented assignment satisfies the complemented clause, so all objective values are unchanged. The formula $T_H$ is therefore a MAX-HORN-2SAT instance up to this renaming, and the Lean development proves both statements.
Every clause has at most two literals, so each formula is also an instance of MAX~2-SAT, and every guarantee proved below for the Horn class is a guarantee for this class of MAX~2-SAT instances, for cubic graphs and $d$-regular graphs alike. On the class $\{T_{H(G)}\}$ of formulas obtained from the conflict graphs of $d$-regular MAX-CUT instances, Theorem~\ref{thm:cutclass} gives the factor $1-\frac{1-\alpha_\varepsilon}{6d+12}\ge83/84$ for every $d\ge1$. The construction below starts from cubic Maximum Independent Set (MIS), uses four clones for every graph vertex, and connects them by a directed cycle of implications. One clone participates in a unit clause and the other three in the three incident edge clauses. Each clone additionally occurs in its incoming and outgoing implication. Consequently every formula variable occurs exactly three times. The reduction from cubic MIS to MAX-HORN-2SAT composes unit, edge, and consistency clauses; the quantitative contribution is the exact optimum identity and the explicit cardinality bound of the repair, together with their Lean formalization.
The contribution is quantitative rather than merely satisfiability preserving. We prove an explicit repair of every inconsistent clone assignment, obtain an exact affine optimum identity, and transfer the Berman--Fujito cubic-MIS guarantee~\cite{BermanFujito1999} to the Horn class. Cubic MIS is APX-complete~\cite{AlimontiKann2000}, and the same repair proves APX-completeness of the image class.
The central result is Theorem~\ref{thm:ugc-class}: the factor $83/84$ on the three-occurrence class $\{T_{H(G)}\}$ lies above the Unique Games thresholds of Brakensiek--Huang--Zwick and Lewin--Livnat--Zwick, which quantify over all instances.
\section{The cubic three-occurrence construction}\label{sec:cubic}
Let $H=(V,E)$ be a finite simple cubic graph, with $n=|V|$ and $m=|E|$. Thus $m=3n/2$. At each vertex $v$, label its incident edges with the three ports $1,2,3$. Introduce four Boolean variables
\[
z_{v,0},z_{v,1},z_{v,2},z_{v,3}.
\]
The formula $T_H$ contains the following clauses.
\begin{enumerate}
\item For every vertex $v$, one unit clause $(\neg z_{v,0})$.
\item For every edge $e=uv$, whose endpoint ports are $p$ at $u$ and $q$ at $v$, one edge clause $(z_{u,p}\vee z_{v,q})$.
\item For every vertex $v$ and $i\in\mathbb Z/4\mathbb Z$, one consistency implication $(\neg z_{v,i}\vee z_{v,i+1})$.
\end{enumerate}
Every clause has at most one negative literal. Every clone occurs in exactly one source clause, one outgoing implication, and one incoming implication; hence it occurs exactly three times. There are $4n$ variables and $n+m+4n=5n+m$ clauses.
For an assignment $z$, let
\begin{align*}
U(z)&=|\{v:z_{v,0}=0\}|,\\
F(z)&=|\{uv\in E:z_{u,p}=z_{v,q}=0\}|,\\
T(z)&=|\{(v,i):z_{v,i}=1,\ z_{v,i+1}=0\}|.
\end{align*}
These are respectively the satisfied unit clauses, falsified edge clauses, and falsified consistency implications. Therefore
\begin{equation}\label{eq:score}
\operatorname{val}_{T_H}(z)=U(z)+(m-F(z))+(4n-T(z)).
\end{equation}
\section{Constructive repair and exact optimum}
Call a vertex selected when $z_{v,0}=0$, and clean when it is selected and its four-bit cycle has no $1$-to-$0$ transition.
\begin{lemma}[Cycle lemma]\label{lem:cycle}
If $z_{v,0}=0$ and the cycle at $v$ has no $1$-to-$0$ transition, then $z_{v,i}=0$ for all $i$.
\end{lemma}
\begin{proof}
A nonconstant cyclic binary word has both a $0$-to-$1$ and a $1$-to-$0$ transition. Since the word begins with zero, absence of a $1$-to-$0$ transition forces it to be constantly zero.
\end{proof}
Let $C(z)$ be the clean selected vertices. Delete the left endpoint of every edge induced by $C(z)$, ignoring repetitions, and call the remaining set $R(z)$.
\begin{lemma}[Quantitative repair]\label{lem:repair}
$R(z)$ is independent and
\[
|R(z)|\ge \operatorname{val}_{T_H}(z)-(4n+m).
\]
The set is computable in polynomial time.
\end{lemma}
\begin{proof}
Every edge induced by $C(z)$ loses an endpoint, so $R(z)$ is independent. At most $T(z)$ selected vertices are nonclean, because each such vertex contributes at least one transition. By Lemma~\ref{lem:cycle}, both endpoint clones of every edge induced by $C(z)$ are zero; hence every such edge is counted by $F(z)$. The number of distinct removed endpoints is at most the number of edges induced by $C(z)$, which is at most $F(z)$. Thus
\[
|R(z)|\ge U(z)-T(z)-F(z).
\]
Substituting \eqref{eq:score} gives the claimed inequality.
\end{proof}
\begin{theorem}[Exact identity]\label{thm:identity}
For every simple cubic graph $H$,
\[
\operatorname{OPT}(T_H)=4n+m+\alpha(H).
\]
Moreover, an assignment of value $A$ decodes in polynomial time to an independent set of size at least $A-(4n+m)$.
\end{theorem}
\begin{proof}
The upper bound follows from Lemma~\ref{lem:repair}. For the lower bound, let $I$ be an independent set and set all four clones of $v$ to zero exactly when $v\in I$. Every consistency implication and every edge clause is satisfied, while precisely $|I|$ unit clauses are satisfied. The value is $4n+m+|I|$. Taking $I$ maximum proves equality.
\end{proof}
\section{The $83/84$ approximation}
For a simple cubic graph put $B=4n+m=11n/2$. Every independent set sends all its incident edges to its complement, so $\alpha\le n/2$. If $|I|\ge6\alpha/7$, the ratio $(B+|I|)/(B+\alpha)$ is at least $(11+6/7)/12=83/84$. We give an endpoint argument that only requires the published cubic MIS factor arbitrarily close to $7/6$~\cite{BermanFujito1999}.
\begin{theorem}\label{thm:ratio}
On the class $\{T_H:H\text{ is simple and cubic}\}$ there is a polynomial-time $83/84$ approximation.
\end{theorem}
\begin{proof}
Return the better lift of two independent sets. First use Berman--Fujito with fixed parameters giving solution fraction $r=6/7-10^{-6}$. If $2\alpha/n\le9999/10000$, then $(83-84r)\alpha\le(11+84\cdot10^{-6})(9999/20000)n \le 11n/2=B$, which gives the required ratio.
For the second set, apply the near-satisfiable Max Cut guarantee of Trevisan~\cite{Trevisan}. For fixed $\epsilon=10^{-4}$, a graph having a cut of fraction at least $1-\epsilon$ admits a polynomially computable cut of fraction at least $0.95$ for all sufficiently large $n$: the stated bound is $1-4\sqrt\epsilon+8\epsilon-o(1)>0.95$. Fixed small orders may be solved exhaustively, including the empty graph. The independent set of size $\alpha$ and its complement cut $3\alpha$ edges, a fraction $2\alpha/n$. Thus the cut guarantee applies in the remaining case. From each side delete one endpoint of each internal edge, taking the union of the deletions on that side. The larger remaining independent set has size at least $(n-0.05m)/2=37n/80>3n/7 \ge6\alpha/7$. The endpoint calculation above now applies. The algorithm runs both branches without knowing $\alpha$.
\end{proof}
\subsection{Generalization to degree $d\ge 3$}
The same construction and ratio transfer extend to every fixed degree $d\ge 3$. Let $H=(V,E)$ be simple and $d$-regular, so $m=dn/2$. At each vertex introduce $d+1$ Boolean clones together with one unit clause, $d$ edge clauses (via a port labelling of the incident edges), and a directed cycle of $d+1$ consistency implications. Every clone still occurs in exactly three places. The score identity becomes
\[
\operatorname{OPT}(T_H)=(d+1)n+m+\alpha(H).
\]
The quantitative repair of Lemma~\ref{lem:repair} is independent of the cycle length, and the same deletion argument yields an independent set of size at least $A-((d+1)n+m)$. Write
\[
B=(d+1)n+m=n\Bigl(\tfrac{3d}{2}+1\Bigr),\qquad
\beta=\tfrac{3d}{2}+1.
\]
Any independent set cuts exactly $d\alpha$ edges, hence $\alpha\le n/2$. Consequently any independent set of size at least $6\alpha/7$ already produces the ratio
\[
\frac{B+|I|}{B+\alpha}
\ge\frac{\beta+3/7}{\beta+1/2}.
\]
The right-hand side is strictly increasing in $\beta$ for $\beta>0$ and equals $83/84$ when $d=3$. Therefore the ratio is at least $83/84$ for every $d\ge 3$.
\begin{proposition}\label{prop:degd}
Let $H$ be simple and $d$-regular, $d\ge3$, and $B=(d+1)n+m$.
\begin{enumerate}
\item The assignment with every clone true has value $B$. It is a $\frac{3d+2}{3d+3}$-approximation on $\{T_H\}$.
\item If $d\ge27$, then $\frac{3d+2}{3d+3}\ge\frac{83}{84}$.
\item If $d\ge22$ and an independent set $I$ satisfies $5\alpha\le(d+3)|I|$, the Berman--Fujito guarantee~\cite{BermanFujito1999}, its lift has value at least $\frac{83}{84}\operatorname{OPT}(T_H)$.
\item If an independent set $I$ satisfies $7|I|\ge6\alpha$, its lift has value at least $\frac{21d+20}{21d+21}\operatorname{OPT}(T_H)$.
\item If $d=21$, there is a polynomial-time $83/84$ approximation on $\{T_H:H\text{ is simple and $21$-regular}\}$.
\item If an independent set $I$ satisfies $28|I|+d\alpha\ge27\alpha$, that is $|I|\ge\rho(d)\alpha$ with $\rho(d)=\frac{27-d}{28}$, its lift has value at least $\frac{83}{84}\operatorname{OPT}(T_H)$.
\item For every fixed $\varepsilon>0$ there is a polynomial-time algorithm on $\{T_H:H\text{ is simple and $d$-regular}\}$ whose value is at least $\Bigl(1-\frac{d-2+\varepsilon}{(3d+3)(d+3+\varepsilon)}\Bigr)\operatorname{OPT}(T_H)$. This exceeds $\frac{3d+2}{3d+3}$ for every $d\ge3$, and the limit $\varepsilon\to0$ is $1-\frac{d-2}{3(d+1)(d+3)}$.
\end{enumerate}
\end{proposition}
\begin{proof}
With every clone true, the $n$ unit clauses fail and the $m$ edge clauses and all $(d+1)n$ consistency implications hold, so the value is $B$. Since $\alpha\le n/2$ and $2B=(3d+2)n$, we get $B\ge(3d+2)\alpha$, which gives the first claim. For the second, $3d+2\ge83$. For the Berman--Fujito claim, multiply the required inequality $B+84|I|\ge83\alpha$ by $d+3$ and use $(d+3)|I|\ge5\alpha$, $(d+3)B\ge(d+3)(3d+2)\alpha$ and $3d^2-72d+177\ge0$ for $d\ge22$. The fourth claim is the endpoint computation above: using $7|I|\ge6\alpha$, $(21d+21)(B+|I|)-(21d+20)(B+\alpha)\ge B-(3d+2)\alpha\ge0$.
For $d=21$ put $a=\alpha/n$, so $B=32.5\,n$ and $m=10.5\,n$. The algorithm returns the better lift of two independent sets, without knowing $a$. The first is a Berman--Fujito set with $\varepsilon=1/50$, which has size at least $50\alpha/241$ since the guarantee for degree $21$ is $24/5+\varepsilon=241/50$. The second is obtained from a cut: solve the Goemans--Williamson semidefinite relaxation of Max Cut on $H$ to additive error $10^{-3}$, round it, and derandomize the rounding by the method of Mahajan and Ramesh~\cite{GoemansWilliamson1995,MahajanRamesh1999}; then delete one endpoint of every edge inside each side, and keep the larger side. If $a\le 99/200$, the first set suffices, because
\[
84\cdot\frac{50}{241}\alpha\ge 83\alpha-B
\quad\Longleftrightarrow\quad
\frac{15803}{241}\,a\le 32.5,
\]
and $\frac{15803}{241}\cdot\frac{99}{200}=\frac{1564497}{48200}<32.5$. If $a>99/200$, the independent set of size $\alpha$ cuts $21\alpha$ of the $m$ edges, a fraction $2a>0.99$, so the relaxation has value at least $c=0.989$. By the analysis of Goemans and Williamson, rounding yields expected cut fraction at least $\arccos(1-2c)/\pi$ whenever $c\ge0.845$; this holds because $\arccos(1-2x)/\pi$ is convex on $[1/2,1]$ and its convex minorant coincides with it for $x\ge0.845$. Here $\arccos(-0.978)/\pi>0.93$. Deleting one endpoint of each uncut edge leaves an independent set of size at least $(n-0.07m)/2=(1-0.735)n/2>0.12\,n$ on the larger side. Then $B+84|I|\ge32.5\,n+10.08\,n>41.5\,n\ge83\alpha$, using $\alpha\le n/2$. In both cases $B+84|I|\ge83\alpha$, which is the required ratio.
For the sixth claim, $B\ge(3d+2)\alpha$ by the argument of the first claim, and the hypothesis gives $84|I|\ge3(27-d)\alpha$. Adding the two gives $B+84|I|\ge83\alpha$. For the seventh claim, the Berman--Fujito set $I$ with $5\alpha\le(d+3+\varepsilon)|I|$ satisfies $\alpha-|I|\le\frac{d-2+\varepsilon}{d+3+\varepsilon}\alpha$, and $B+\alpha\ge(3d+3)\alpha$, so $\frac{B+|I|}{B+\alpha}=1-\frac{\alpha-|I|}{B+\alpha}\ge1-\frac{d-2+\varepsilon}{(3d+3)(d+3+\varepsilon)}$; the loss is below $\frac1{3d+3}$ because $d-2+\varepsilon<d+3+\varepsilon$.
\end{proof}
The threshold $\rho(d)=(27-d)/28$ equals $6/7$ for $d=3$. Selected values:
\begin{center}
\begin{tabular}{c|cccccccc}
$d$ & 3 & 4 & 5 & 6 & 8 & 10 & 12 & 16 \\
\hline
$\rho(d)$ & $6/7$ & $23/28$ & $11/14$ & $3/4$ & $19/28$ & $17/28$ & $15/28$ & $11/28$
\end{tabular}
\end{center}
The value $83/84$ for $d=3$ uses the Berman--Fujito branch and the near-bipartite branch of Theorem~\ref{thm:ratio}. The $6/7$ hypothesis in the fourth claim is the hypothesis whose cubic instance is supplied by that theorem.
For $\varepsilon\to0$ the ratio of the seventh claim is $1-\frac{d-2}{3(d+1)(d+3)}$. Selected values:
\begin{center}
\begin{tabular}{c|cccc}
$d$ & 4 & 5 & 6 & 20 \\
\hline
$1-\frac{d-2}{3(d+1)(d+3)}$ & $\frac{103}{105}$ & $\frac{47}{48}$ & $\frac{185}{189}$ & $\frac{159}{161}$ \\
$\frac{3d+2}{3d+3}$ & $\frac{14}{15}$ & $\frac{17}{18}$ & $\frac{20}{21}$ & $\frac{62}{63}$
\end{tabular}
\end{center}
Explicitly, the endpoint ratio obtained from a $6/7$-approximate independent set is
\[
\frac{21d+20}{21d+21}.
\]
Selected values:
\begin{center}
\begin{tabular}{c|ccccccc}
$d$ & 3 & 4 & 5 & 6 & 8 & 10 & 12 \\
\hline
ratio & $83/84$ & $104/105$ & $125/126$ & $146/147$ & $188/189$ & $230/231$ & $272/273$ \\
& $\approx0.9881$ & $\approx0.9905$ & $\approx0.9921$ & $\approx0.9932$ & $\approx0.9947$ & $\approx0.9957$ & $\approx0.9963$
\end{tabular}
\end{center}
As $d\to\infty$ both $\frac{3d+2}{3d+3}$ and $\frac{21d+20}{21d+21}$ tend to $1$.
\section{APX-completeness}
\begin{theorem}\label{thm:apx}
The cubic three-occurrence Horn class is APX-complete under $L$-reductions. The reduction from cubic MIS has constants $23$ and $1$.
\end{theorem}
\begin{proof}
A greedy maximal independent set in a graph of maximum degree three has size at least $n/4$, so $n\le4\alpha(H)$. Since $m=3n/2$,
\[
\operatorname{OPT}(T_H)=\frac{11n}{2}+\alpha(H)
\le22\alpha(H)+\alpha(H)=23\alpha(H).
\]
For an assignment of value $A$, Lemma~\ref{lem:repair} returns $R$ with $|R|\ge A-(4n+m)$. Hence
\[
\alpha(H)-|R|
\le (4n+m+\alpha(H))-A
=\operatorname{OPT}(T_H)-A.
\]
These are the two $L$-reduction inequalities. Cubic MIS is APX-complete and Theorem~\ref{thm:ratio} places the image class in APX.
\end{proof}
\subsection{Hardness for larger degree}
The same argument yields APX-completeness for every fixed $d\ge 3$. Maximum Independent Set remains APX-complete on simple $d$-regular graphs (and on graphs of maximum degree $d$) for each fixed $d\ge 3$ \cite{AlimontiKann2000,BermanFujito1999}. The construction of $T_H$ from a $d$-regular graph is an $L$-reduction whose constants depend only on $d$.
A greedy maximal independent set in a degree-$d$ graph has size at least $n/(d+1)$, so $n\le(d+1)\alpha$. With $m=dn/2$ one obtains
\[
\operatorname{OPT}(T_H)=(d+1)n+m+\alpha
\le(d+1)^2\Bigl(\tfrac{3d}{2}+1\Bigr)\alpha+\alpha.
\]
Thus the first $L$-reduction constant is $O(d^2)$. The second constant remains exactly $1$ by the quantitative repair: every assignment of value $A$ produces an independent set $R$ satisfying $\alpha-|R|\le\operatorname{OPT}(T_H)-A$. Consequently the image class
\[
\{T_H:H\text{ is simple and $d$-regular}\}
\]
is APX-complete for every fixed $d\ge 3$.
By Proposition~\ref{prop:degd}(1) the image class is approximable to the factor $1-\frac1{3d+3}$ for every $d$, so it lies in APX and is APX-complete.
Under the Unique Games Conjecture the unrestricted weighted MAX-HORN-2SAT problem is hard to approximate above the constant $\alpha^*\approx0.94616$ of Brakensiek--Huang--Zwick. The factor $\frac{3d+2}{3d+3}$ exceeds $\alpha^*$ for every $d\ge6$, and Theorem~\ref{thm:ratio} gives $83/84>\alpha^*$ for $d=3$. Both statements concern bounded-occurrence classes and are consistent with the bounded-occurrence improvements of Ghoshal et al.~\cite{GHLMM2026} and with the Unique Games Conjecture.
\section{A signed graph for general weighted inputs}\label{sec:signed}
\paragraph{Unary terminals.} For the path $a$--$w$--$b$, the terminal states $00$, $01$, $10$, and $11$ give maximum total cardinalities $1,1,1,2$. It realizes a one-unit bonus for $11$ when another constraint synchronizes the ends. Both mixed states extend to independent sets. The following signed construction instead gives a unit clause one reward vertex attached to its falsifying port.
\begin{theorem}\label{thm:signed}
For any weighted source $F$ of width at most two, let $D$ be its number of literal occurrences, $W$ its total nonnegative clause weight, and $M=W+1$. One can construct a simple vertex-weighted graph $J$ of maximum degree three such that $\alpha_w(J)=DM+\operatorname{OPT}(F)$. Every independent set of weight $A$ decodes in polynomial time to an assignment of $F$ of score at least $A-DM$.
\end{theorem}
\begin{proof}
At occurrence $i$ create adjacent ports $T_i,F_i$ of weight $M$. Join $T_i$ to $F_{i+1}$ around each variable cycle. For a singleton cycle the pair edge already supplies this conflict; insert it just once. Discard variables with no occurrences. For each literal occurrence create a reward vertex of weight equal to its clause weight and join it to the port that falsifies that literal. Join the two reward vertices of each binary clause. Every port has at most two ring neighbors and one reward neighbor. Distinct literal occurrences have distinct vertices even for repeated literals. The graph has $3D$ vertices, linear combinatorial size and polynomial bit complexity with binary rational weights.
At most one port per pair and one reward per clause can be selected. If fewer than $D$ ports are selected, the weight is at most $(D-1)M+W<DM$. Any constant-per-variable encoding already reaches $DM$. Otherwise each pair has exactly one selected port. Choosing $T_i$ forces $T_{i+1}$, so all pairs of a variable choose the same side. Each selected reward then witnesses a true literal. Reading the ports gives score at least $A-DM$. Conversely encode any assignment and select one true witness per satisfied clause. The weight is exactly $DM$ plus its score. These operations prove both claims, including for empty clauses.
\end{proof}
This construction preserves multiplicative guarantees on the excess $A-DM$. A guarantee on the total weight alone can be almost perfect when the decoded score is zero: for $k$ positive unit clauses on one variable, the all-false port selection has fraction $(k+1)/(k+2)$ and decodes to zero. Already $k=82$ gives $83/84$. Thus neither the signed repair nor a polynomial algorithm for unweighted cubic MIS supplies the unrestricted weighted ratio.
\section{Unique Games thresholds, MAX~2-SAT and MAX-CUT}
\label{sec:ugc}
Let $\alpha^*$ denote the constant of Brakensiek, Huang and Zwick for weighted MAX~$\{1,2\}$-HORN-SAT, approximately $0.94615981$. Their Theorem~5.12 gives UG-hardness at $\alpha^*+\varepsilon$ for every $\varepsilon>0$~\cite{BrakensiekHuangZwick2024}. In particular, $19/20=0.95>\alpha^*$. A deterministic polynomial-time algorithm returning value at least $19\operatorname{OPT}(F)/20$ on \emph{every} weighted source $F$ would solve the UG-hard gap problem in polynomial time and, assuming UGC, imply $\mathrm P=\mathrm{NP}$. The worst-case bound in Theorem~\ref{thm:ratio} concerns the cubic family.
The constant $\beta^\star\approx0.94016567$ of Lewin, Livnat and Zwick is the UG-hard threshold for MAX~2-SAT~\cite{BrakensiekHuangZwick2024}. Both UG-hardness statements quantify over all instances. They give no information about the class in which every variable has exactly three literal occurrences. For that class Theorem~\ref{thm:ratio} gives the worst-case ratio $83/84$, and
\[
\frac{83}{84}>0.94615981>0.94016567 .
\]
Ghoshal, Huang, Lee, Makarychev and Makarychev prove that MAX~2-SAT instances in which every variable appears in at most $d$ constraints admit a $\beta^\star+\tilde\Omega(1/d^{2})$-approximation~\cite{GHLMM2026}. Thus a guarantee above $\beta^\star$ on a bounded-occurrence class is established for every fixed $d$ and is consistent with the Unique Games Conjecture. A guarantee above $\alpha^*$ or $\beta^\star$ on \emph{all} instances would not be. For MAX~2-SAT instances in which every variable appears in at most two clauses, the problem can be solved exactly in linear time, while already three occurrences per variable give an NP-complete problem~\cite{RamanRavikumarRao1998}. Thus the occurrence bound $d=3$ is the first at which the problem is NP-hard, and for every fixed $d$ the guarantee of~\cite{GHLMM2026} exceeds $\beta^\star$.
The endpoint guarantee obtained from a $6/7$-approximate independent set on a $d$-regular graph is
\[
r(d)=\frac{21d+20}{21d+21},
\]
which is strictly increasing for $d\ge3$ since
\[
r(d+1)-r(d)=\frac{21}{(21d+21)(21d+42)}>0 .
\]
For $d=3$ the independent-set hypothesis is supplied by Theorem~\ref{thm:ratio}, and $r(3)=83/84$ exceeds $\beta^\star$ by $\approx0.04793$. For every $d\ge6$ Proposition~\ref{prop:degd}(1) gives the unconditional factor $\frac{3d+2}{3d+3}>0.95>\alpha^*>\beta^\star$ on the $d$-regular family. For $d\ge3$ one has $\frac{d-2}{3(d+1)(d+3)}\le\frac4{189}<0.022$, because $63(d-2)\le4(d+1)(d+3)$ is equivalent to $(4d-23)(d-6)\ge0$, which holds for every integer $d$. Hence for every $d\ge4$ and every sufficiently small $\varepsilon>0$ the factor of Proposition~\ref{prop:degd}(7) exceeds $0.97>\alpha^*>\beta^\star$ on the $d$-regular family, and together with $r(3)=83/84$ every class $\{T_H:H\text{ is simple and $d$-regular}\}$ with $d\ge3$ has a polynomial-time factor above $\alpha^*$. These are statements about three-occurrence classes and do not concern all instances. The additive gain $\tilde\Omega(1/d^{2})$ of~\cite{GHLMM2026} is $o(1)$ in $d$. By Proposition~\ref{prop:degd}(6), for each $d$ with $4\le d\le20$ an independent-set algorithm of ratio $\rho(d)=(27-d)/28$ on $d$-regular graphs would give $83/84$ on the $d$-regular family, and $83/84>\alpha^*>\beta^\star$. The thresholds of Brakensiek, Huang and Zwick quantify over all instances. OpenAI released on 6~October 2026 a manuscript claiming a proof of the Unique Games Conjecture, together with a Lean formalization; this claim has not been peer reviewed~\cite{OpenAIMath2026}. If it is confirmed, the UG-hardness statements above become NP-hardness statements.
\subsection{MAX-CUT as MAX 2-SAT}
\begin{proposition}\label{prop:cut2sat}
Let $G=(V,E)$ be a graph with $n$ vertices and $m$ edges. Let $F_G$ be the formula with one variable $x_v$ for every vertex and, for every edge $uv$, the two clauses $(x_u\lor x_v)$ and $(\lnot x_u\lor\lnot x_v)$. Then:
\begin{enumerate}
\item an assignment $z$ satisfies exactly $m+\operatorname{cut}(S_z)$ clauses, where $S_z=\{v:z(x_v)=\text{true}\}$ and $\operatorname{cut}(S)$ is the number of edges with exactly one endpoint in $S$; hence $\operatorname{OPT}(F_G)=m+\operatorname{maxcut}(G)$;
\item the variable $x_v$ occurs in exactly $2\deg(v)$ clauses;
\item if an assignment satisfies at least $\rho\operatorname{OPT}(F_G)$ clauses, then the cut $S_z$ has at least $\rho\,(m+\operatorname{maxcut}(G))-m$ edges.
\end{enumerate}
\end{proposition}
\begin{proof}
For an edge $uv$ whose endpoints receive different values both clauses hold. If both endpoints are true only $(x_u\lor x_v)$ holds, and if both are false only $(\lnot x_u\lor\lnot x_v)$ holds. Summing over the edges gives the first claim, and the second and third follow. \end{proof}
The clause $(\lnot x_u\lor\lnot x_v)$ has two negative literals, so $F_G$ is an instance of MAX~2-SAT that does not belong to the Horn class; no guarantee of this paper is claimed for it.
\begin{proposition}\label{prop:cutconflict}
Let $G$ be a simple $d$-regular graph with $n$ vertices and $m=dn/2$ edges. Let $H$ be the conflict graph of $F_G$: its $N=4m$ vertices are the literal occurrences of $F_G$, and two vertices are adjacent if they lie in the same clause or are complementary literals. Then $H$ is simple and $(d+1)$-regular, and
\[
\alpha(H)=\operatorname{OPT}(F_G)=m+\operatorname{maxcut}(G).
\]
Consequently, with $D=d+1$, $M=DN/2$ and $B=(D+1)N+M=2m(3d+5)$, the formula $T_H$ satisfies $\operatorname{OPT}(T_H)=B+m+\operatorname{maxcut}(G)$, and an assignment of $T_H$ of value $A$ yields in polynomial time a cut of $G$ with at least $A-B-m$ edges.
\end{proposition}
\begin{proof}
Every literal occurrence has exactly one neighbour in its clause. A positive occurrence of $x_u$ is adjacent to the $d$ negative occurrences of $x_u$, and symmetrically; an in-clause neighbour is a literal of another variable, so no edge is repeated. Hence $H$ is simple and $(d+1)$-regular. An independent set contains at most one literal of each clause and no two complementary literals, so it is a consistent partial assignment satisfying as many clauses as the set has vertices; conversely, an assignment and one satisfied literal per satisfied clause give an independent set. Thus $\alpha(H)=\operatorname{OPT}(F_G)$, which equals $m+\operatorname{maxcut}(G)$ by Proposition~\ref{prop:cut2sat}. The identity for $T_H$ is the degree-$D$ exact identity. By the quantitative repair, an assignment of value $A$ gives an independent set of size at least $A-B$, hence an assignment of $F_G$ satisfying at least $A-B$ clauses, hence a cut with at least $A-B-m$ edges by Proposition~\ref{prop:cut2sat}.
\end{proof}
\begin{corollary}\label{cor:cutratio}
Let $\operatorname{maxcut}(G)=\gamma m$ and $0\le\rho\le1$.
\begin{enumerate}
\item If an assignment of $F_G$ satisfies at least $\rho\operatorname{OPT}(F_G)$ clauses, its cut has at least $\bigl(\rho-\frac{1-\rho}{\gamma}\bigr)\operatorname{maxcut}(G)$ edges.
\item If $G$ is $d$-regular and an assignment of $T_H$ has value at least $\rho\operatorname{OPT}(T_H)$, the cut obtained from it has at least $\bigl(\rho-\frac{(1-\rho)(6d+11)}{\gamma}\bigr)\operatorname{maxcut}(G)$ edges.
\end{enumerate}
\end{corollary}
\begin{proof}
By Proposition~\ref{prop:cut2sat}, the cut has at least $\rho(m+\gamma m)-m$ edges; dividing by $\gamma m$ gives the first claim. By Proposition~\ref{prop:cutconflict} it has at least $\rho(B+m+\gamma m)-B-m$ edges with $B+m=(6d+11)m$, which gives the second.
\end{proof}
\begin{theorem}\label{thm:cutclass}
For every $d\ge1$ there is a polynomial-time algorithm that, given a simple $d$-regular graph $G$, returns an assignment of $T_{H(G)}$ of value at least
\[
\Bigl(1-\frac{1-\alpha_\varepsilon}{6d+12}\Bigr)\operatorname{OPT}(T_{H(G)})\ \ge\ \frac{83}{84}\operatorname{OPT}(T_{H(G)}),
\]
where $H(G)$ is the conflict graph of Proposition~\ref{prop:cutconflict} and $\alpha_\varepsilon\approx0.87856$ Here $\alpha_\varepsilon=\alpha_{\mathrm{GW}}-10^{-5}\ge0.8785$, where $\alpha_{\mathrm{GW}}\approx0.87856$ is the Goemans--Williamson constant.
\end{theorem}
\begin{proof}
Let $c=\operatorname{maxcut}(G)$. Solving the Goemans--Williamson relaxation to additive error $10^{-5}$ and rounding with derandomization~\cite{GoemansWilliamson1995,MahajanRamesh1999} returns in polynomial time a set $S$ with $\operatorname{cut}(S)\ge\alpha_\varepsilon\,c$ (the relaxation error $10^{-5}$ is absorbed in $\alpha_\varepsilon$). Set $x_v$ true exactly for $v\in S$ and choose one true literal in every satisfied clause of $F_G$. These $m+\operatorname{cut}(S)$ literal occurrences are pairwise non-adjacent in $H(G)$, because they lie in distinct clauses and no two are complementary. So they form an independent set $I$ with $|I|=m+\operatorname{cut}(S)$, while $\alpha(H(G))=m+c$ by Proposition~\ref{prop:cutconflict}. Since $c\le m$,
\[
|I|\ge m+\alpha_\varepsilon\,c\ge\frac{1+\alpha_\varepsilon}{2}\,(m+c)=\frac{1+\alpha_\varepsilon}{2}\,\alpha(H(G)),
\]
because $(1-\alpha_\varepsilon)(m-c)\ge0$. Let $D=d+1$, so $H(G)$ is $D$-regular and $B=(D+1)N+M\ge(3D+2)\alpha(H(G))$ by $\alpha(H(G))\le N/2$. The lift of $I$ has value $B+|I|$, hence
\[
\frac{B+|I|}{B+\alpha(H(G))}\ge1-\frac{\frac{1-\alpha_\varepsilon}{2}\,\alpha(H(G))}{B+\alpha(H(G))}\ge1-\frac{1-\alpha_\varepsilon}{6d+12}.
\]
Finally $\frac{1-\alpha_\varepsilon}{6d+12}\le\frac{1-\alpha_\varepsilon}{18}<\frac{1}{84}$ for $d\ge1$, since $1-\alpha_\varepsilon\le0.1215<18/84\approx0.214$.
\end{proof}
\begin{proposition}\label{prop:bounded-cut}
Let $c_d\in(0,1]$ and suppose a polynomial-time algorithm returns on every graph of maximum degree $d$ a cut with at least $c_d\operatorname{maxcut}(G)$ edges; if the algorithm is randomized and the guarantee is for the expected cut, every statement below is for the expected value. Then on $\{T_{H(G)}:G\text{ is }d\text{-regular}\}$ a polynomial-time algorithm has the ratio at least $1-\frac{1-c_d}{6d+12}$.
\end{proposition}
\begin{proof}
As in the proof of Theorem~\ref{thm:cutclass} with $0.8785$ replaced by $c_d$: $|I|\ge\frac{1+c_d}{2}\alpha(H(G))$ and $B+\alpha(H(G))\ge(3D+3)\alpha(H(G))$ with $D=d+1$, so the loss is at most $\frac{1-c_d}{2(3d+6)}$.
\end{proof}

For graphs of maximum degree $d$, Feige, Karpinski and Langberg give a guarantee $\alpha_{\mathrm{GW}}+\epsilon_d$ with $\epsilon_d>0$, and $0.921$ for $d=3$~\cite{FeigeKarpinskiLangberg2002}; Halperin, Livnat and Zwick give an algorithm whose expected approximation ratio on graphs of maximum degree $3$ is at least $0.9326$~\cite{HalperinLivnatZwick2004}; Hsieh and Kothari give $\alpha_{\mathrm{GW}}+\tilde\Omega(1/d^2)$~\cite{HsiehKothari2023}. With $c_3=0.9326$ as an expected guarantee, Proposition~\ref{prop:bounded-cut} gives the expected ratio $1-\frac{0.0674}{30}>0.9977$ on $\{T_{H(G)}:G\text{ cubic}\}$. These bounded-degree guarantees exceed $\alpha_{\mathrm{GW}}$ on a class of graphs and are consistent with the threshold of~\cite{KKMO2007}, which quantifies over all graphs.

\begin{remark}\label{rem:cutclass-range}
Theorem~\ref{thm:cutclass} applies to every degree $D=d+1$ with $4\le D\le20$, that is $3\le d\le19$, on the class $\{T_{H(G)}:G\text{ is }(D-1)\text{-regular}\}$. In that range $|I|\ge\frac{1+\alpha_\varepsilon}{2}\,\alpha(H(G))$ exceeds $\rho(D)=\frac{27-D}{28}\le\frac{23}{28}$, so Proposition~\ref{prop:degd}(6) gives the ratio $83/84$ without any assumption on an independent-set algorithm. Since $\frac{1+\alpha_\varepsilon}{2}>6/7$, Proposition~\ref{prop:degd}(4) gives the ratio $\frac{21D+20}{21D+21}$, which equals $\frac{104}{105}$ for $D=4$ and $\frac{440}{441}$ for $D=20$. For these $G$ the constant of Corollary~\ref{cor:cutratio} is $B+m=(6D+5)m$; for $D=4$ it equals $29m$, and an assignment of $T_{H(G)}$ of ratio $\frac{104}{105}$ yields a cut with at least $\frac{104}{105}\operatorname{maxcut}(G)-\frac{29}{105}m$ edges, which is $\frac{49}{105}m$ when $\operatorname{maxcut}(G)=\frac34m$. The statement concerns the conflict-graph class; for arbitrary $D$-regular graphs $H$ with $4\le D\le20$ the ratio $83/84$ is given by Proposition~\ref{prop:degd}(6) under the hypothesis stated there.
\end{remark}
\subsection{Consequences of the MAX-CUT reduction for the Unique Games Conjecture}
Khot, Kindler, Mossel and O'Donnell show that, assuming the Unique Games Conjecture, MAX-CUT cannot be approximated in polynomial time within any constant above $\alpha_{\mathrm{GW}}=\min_{0<x\le1}\frac{\arccos(1-2x)}{\pi x}\approx0.87856$~\cite{KKMO2007}, the guarantee of the Goemans--Williamson algorithm. This threshold quantifies over all graphs. Theorem~\ref{thm:cutclass} concerns the formulas $T_{H(G)}$ obtained from MAX-CUT instances $G$ through Proposition~\ref{prop:cut2sat} and Proposition~\ref{prop:cutconflict}. Every variable of such a formula occurs exactly three times, and the guarantee is $83/84$ for $d$-regular $G$, above both $\alpha^*$ and $\beta^\star$. These are statements about a three-occurrence class and are therefore consistent with the hardness statements of this section.
\begin{theorem}[Ratio above the Unique Games thresholds]\label{thm:ugc-class}
For every $d\ge1$ there is a polynomial-time algorithm that, given a simple $d$-regular graph $G$, returns an assignment of the three-occurrence Horn formula $T_{H(G)}$ of value at least $\frac{83}{84}\operatorname{OPT}(T_{H(G)})$, and
\[
\frac{83}{84}>\alpha^*\approx0.94616>\beta^\star\approx0.94017 .
\]
Under the Unique Games Conjecture no polynomial-time algorithm has a factor above $\alpha^*+\varepsilon$ on all instances of weighted MAX~$\{1,2\}$-HORN-SAT, or above $\beta^\star+\varepsilon$ on all instances of MAX~2-SAT, for any $\varepsilon>0$. The class $\{T_{H(G)}\}$ is a class of Horn formulas, hence of MAX~2-SAT instances, in which every variable has exactly three literal occurrences, and it is consistent with the conjecture.
\end{theorem}
\begin{proof}
The algorithm and the factor $83/84$ are Theorem~\ref{thm:cutclass}. The numerical inequalities and the hardness statements are those quoted in the first paragraph of Section~\ref{sec:ugc} and~\cite{BrakensiekHuangZwick2024}.
\end{proof}
Take the ratio $\rho_d=1-\frac{1-\alpha_\varepsilon}{6d+12}$ of Theorem~\ref{thm:cutclass}. Then the cut fraction guaranteed by Corollary~\ref{cor:cutratio}(2) is
\[
1-\frac{1-\alpha_\varepsilon}{6d+12}\Bigl(1+\frac{6d+11}{\gamma}\Bigr)\le1-\frac{1-\alpha_\varepsilon}{6d+12}\,(6d+12)=\alpha_\varepsilon\qquad\bigl(\gamma=\operatorname{maxcut}(G)/m\in(0,1]\bigr),
\]
with equality exactly for $\gamma=1$. This is the worst case over $\gamma$ of the corollary and is not a property of the algorithm. The algorithm does not obtain its cut from $\rho_d$. It first computes $S$ with $\operatorname{cut}(S)\ge\alpha_\varepsilon\operatorname{maxcut}(G)$ and then lifts the independent set $I$ of size $m+\operatorname{cut}(S)$ of Theorem~\ref{thm:cutclass}. The cut recovered from the lift is $S$ itself, so the MAX-CUT guarantee is the one of~\cite{GoemansWilliamson1995,MahajanRamesh1999}, and the reduction transfers it to a guarantee on $T_{H(G)}$ that exceeds $\alpha^*$; the MAX-CUT guarantee is not improved. A bounded-degree cut guarantee $c_d$ enters only through Proposition~\ref{prop:bounded-cut}.
Lean proves the clause-count identity $m+\operatorname{cut}(S)$ and the arithmetic of Corollary~\ref{cor:cutratio} (\texttt{CutTransfer}), the lifting step with $18785\,\alpha(H)\le20000\,|I|$, which follows from $|I|\ge\frac{1+\alpha_\varepsilon}{2}\alpha(H)$ (\texttt{CutLift}), and the division-free form of Proposition~\ref{prop:bounded-cut} (\texttt{CitedResults}). The Goemans--Williamson guarantee, the bounded-degree cut guarantees and the Khot--Kindler--Mossel--O'Donnell threshold are taken from the literature (\texttt{CitedResults}) and are not formalized.
\section{Lean verification}
The module \texttt{PalomaMax2SAT.CubicThreeOccurrence} represents the cubic graph with three ports at every vertex. It formalizes the transition count, false edge clauses, explicit deletion repair, independence of the repaired set, its cardinality bound, the matching assignment from an independent set, and the exact optimum identity. It also proves the division-free inequality
\[
84(4n+m+|I|)\ge83(4n+m+\alpha)
\]
from $2m=3n$, $2\alpha\le n$, and $7|I|\ge6\alpha$, and the rational identity $11+6/7=(83/84)\cdot12$. The file \texttt{Solution.lean} proves directly the general arithmetic transfer
\[
84\bigl((d+1)n+m+|I|\bigr)\ge83\bigl((d+1)n+m+\alpha\bigr)
\]
valid for every $d\ge3$ under the analogous hypotheses $2m=dn$ and $2\alpha\le n$, and the monotonicity statement that the endpoint ratio \((\beta+3/7)/(\beta+1/2)\) is at least \(83/84\) for all \(\beta\ge11/2\). GitHub Actions builds the module and runs the public theorem comparison.
The module \texttt{PalomaMax2SAT.SemanticHorn} defines $T_H$ explicitly as a family of clauses, each a list of literals over the clone variables, indexed by the vertices (unit clauses), the edges (edge clauses) and the pairs $(v,i)$ (consistency implications). It proves that the number of clauses of $T_H$ satisfied by $z$ equals $U(z)+(m-F(z))+(4n-T(z))$, that the maximum of this number over all assignments is exactly $4n+m+\alpha$, that every clause has at most one negative literal, that every clause becomes a Horn clause when every variable is renamed to its negation and that this renaming preserves satisfaction of each clause, that $T_H$ has $5n+m$ clauses, and that every variable occurs exactly three times. The Berman--Fujito guarantee, the Trevisan near-bipartite branch, the Goemans--Williamson rounding, the APX-completeness of cubic Maximum Independent Set, the discussion of the Unique Games Conjecture are proved in the text or taken from the literature and are not formalized in Lean.
The module \texttt{PalomaMax2SAT.DegreeTransfer} formalizes the arithmetic of Proposition~\ref{prop:degd}. From $2m=dn$ and $2\alpha\le n$, the value $(d+1)n+m$ satisfies $(3d+2)(B+\alpha)\le(3d+3)B$ with $B=(d+1)n+m$. For $d\ge27$ it satisfies $83(B+\alpha)\le84B$. For $d\ge22$ and $5\alpha\le(d+3)s$ one gets $83(B+\alpha)\le84(B+s)$. From $6\alpha\le7s$ one gets $(21d+20)(B+\alpha)\le(21d+21)(B+s)$, and from $27\alpha\le28s+d\alpha$ one gets $83(B+\alpha)\le84(B+s)$. For $d=21$ it proves the two-branch arithmetic: from $2m=21n$, $2\alpha\le n$, $200\alpha\le99n\Rightarrow50\alpha\le241s$ and $99n<200\alpha\Rightarrow3n\le25s$ one gets $83(22n+m+\alpha)\le84(22n+m+s)$. It also proves the two branches of the $d=21$ argument separately, the inequality $\frac{15803}{241}\cdot\frac{99}{200}<\frac{65}{2}$, the fact that $200\alpha>99n$ gives a cut fraction above $0.99$, the step from at most $7m/100$ uncut edges to $3n\le25|I|$, and the tabulated values of $\rho(d)$ and of $(21d+20)/(21d+21)$. The independent-set guarantees are hypotheses; Lean does not formalize the semidefinite rounding. The module \texttt{PalomaMax2SAT.DegreeFormula} defines $T_H$ for every $d$ as an explicit list of Horn clauses over the $d+1$ clones of each vertex. It proves that the number of satisfied clauses equals $U(z)+(m-F(z))+((d+1)n-T(z))$, that the optimum is exactly $(d+1)n+m+\alpha$, that every clause has at most one negative literal and becomes a Horn clause when every variable is renamed to its negation, that $T_H$ has $(d+2)n+m$ clauses and that every variable occurs exactly three times. It derives $2m=dn$ from the port structure and $2\alpha\le n$ for $d>0$, proves that the lift of an independent set $I$ has value $(d+1)n+m+|I|$, and combines these with the arithmetic of \texttt{DegreeTransfer} into statements about $T_H$ itself: the baseline ratio $\frac{3d+2}{3d+3}$, the ratio $83/84$ for $d\ge27$, and the transfers of the Berman--Fujito, $6/7$, threshold and $d=21$ hypotheses.
For every degree $e+2$, \texttt{DegreeTransfer} proves the division-free form of the seventh claim of Proposition~\ref{prop:degd}, with $\varepsilon=1/c$ for a natural number $c$, and \texttt{DegreeFormula} proves it for the lift of an independent set of the formula $T_H$. The arithmetic inequality $63d\le4(d+1)(d+3)+126$ used in the sweep over $d\ge3$ is also proved in \texttt{DegreeTransfer}. The module \texttt{PalomaMax2SAT.CutTransfer} proves Proposition~\ref{prop:cut2sat}(1), that an assignment of $F_G$ satisfies exactly $m+\operatorname{cut}(S)$ clauses, and the arithmetic of the two ratio transfers of Corollary~\ref{cor:cutratio}. The conflict-graph identity of Proposition~\ref{prop:cutconflict} is proved in \texttt{ConflictGraph}. The module \texttt{PalomaMax2SAT.CutLift} proves that an independent set $I$ of $H$ with $18785\,\alpha(H)\le20000\,|I|$ lifts to the ratio $83/84$ on $T_H$ for every degree $D\ge1$, which is the last step of Theorem~\ref{thm:cutclass}; the construction of $I$ from a cut is in \texttt{ConflictGraph} and the cut guarantee is taken from the literature. The module \texttt{PalomaMax2SAT.CitedResults} states the published guarantees used by the $d=21$ argument as explicit hypotheses, namely the existence of a cut with at most $7m/100$ uncut edges~\cite{GoemansWilliamson1995,MahajanRamesh1999} and of a Berman--Fujito independent set~\cite{BermanFujito1999}. It proves that such a cut yields an independent set $I$ with $3n\le25|I|$ and that, under the two hypotheses, some independent set lifts to a $83/84$-approximation of $T_H$ for $d=21$. It also states the Goemans--Williamson cut guarantee and the bounded-degree cut guarantees as numerical hypotheses, from which it proves the independent-set step of Theorem~\ref{thm:cutclass} and the ratio of Proposition~\ref{prop:bounded-cut}. The module \texttt{PalomaMax2SAT.ConflictGraph} proves $\alpha(H(G))=m+\operatorname{maxcut}(G)$ for every graph given by its edge list, and \texttt{PalomaMax2SAT.ConflictDegree} proves that the conflict graph is $(d+1)$-regular when $G$ is $d$-regular, that $x_v$ occurs $2\deg(v)$ times and that $H(G)$ has $4m$ vertices. The module \texttt{PalomaMax2SAT.LReduction} proves that a graph of maximum degree $\Delta$ has an independent set $I$ with $n\le(\Delta+1)|I|$, the first $L$-reduction constant $\operatorname{OPT}(T_H)\le23\,\alpha(H)$ for cubic graphs and $2\operatorname{OPT}(T_H)\le((3d+2)(d+1)+2)\alpha(H)$ for $d$-regular graphs, and the second constant $1$ (Theorem~\ref{thm:apx}); the APX-completeness of cubic Maximum Independent Set is cited. The module \texttt{PalomaMax2SAT.PaperArithmetic} proves the comparison $\frac{83}{84}>0.94615981>0.94016567$ of Theorem~\ref{thm:ugc-class}, the two-branch arithmetic of Theorem~\ref{thm:ratio} with the two branch guarantees as hypotheses, the numerical constants of Theorem~\ref{thm:cutclass} and of Proposition~\ref{prop:bounded-cut}, and the tabulated values. The module \texttt{PalomaMax2SAT.SignedGraph} proves Theorem~\ref{thm:signed}: the independent sets of the signed graph are the pairs of a port set and a reward set satisfying the four edge conditions, their weight is at most $DM+\operatorname{OPT}(F)$, an independent set of weight $A$ decodes to an assignment of score at least $A-DM$, and the bound is attained; the vertex count and the maximum degree of the graph are not formalized. The repository is available at~\cite{VegaPalomaRepository}. The following statements of the paper are not formalized in Lean: the APX-completeness statement of Theorem~\ref{thm:apx} (only the $L$-reduction constants are proved); the vertex count and the maximum degree of the signed graph and the two examples of Section~\ref{sec:signed}; the simplicity of $H(G)$ as a separate lemma; $\operatorname{OPT}(F_G)=m+\operatorname{maxcut}(G)$ and $\operatorname{OPT}(T_{H(G)})=B+m+\operatorname{maxcut}(G)$ as single statements (their ingredients are proved) and Corollary~\ref{cor:cutratio} with real $\gamma$ and $\rho$ (it is proved division-free); the decimal remarks $19/20>\alpha^*$, $0.04793$ and $0.022$ of Section~\ref{sec:ugc}; the Trevisan bound, the Goemans--Williamson bound and the exhaustive solution of small orders in Theorem~\ref{thm:ratio} and Proposition~\ref{prop:degd}(5); the polynomial running time of every constructive statement; the semidefinite rounding, the Berman--Fujito, Trevisan, Feige--Karpinski--Langberg, Halperin--Livnat--Zwick and Hsieh--Kothari guarantees, and the Unique Games Conjecture with its hardness thresholds. The notes in \texttt{doc/formalization.md} list each cited result with its reference.
\begin{thebibliography}{9}
\bibitem{AlimontiKann2000}
P.~Alimonti and V.~Kann, \emph{Some APX-completeness results for cubic graphs}, Theoretical Computer Science 237 (2000), 123--134. \url{https://doi.org/10.1016/S0304-3975(98)00158-3}.
\bibitem{Karp1972}
R.~M. Karp, \emph{Reducibility among combinatorial problems}, in Complexity of Computer Computations, Plenum, 1972, 85--103.
\bibitem{Hastad1999}
J.~H{\aa}stad, \emph{Clique is hard to approximate within $n^{1-\epsilon}$}, Acta Mathematica 182 (1999), 105--142.
\bibitem{Zuckerman2007}
D.~Zuckerman, \emph{Linear degree extractors and the inapproximability of max clique and chromatic number}, Theory of Computing 3 (2007), 103--128.
\bibitem{Khot2002}
S.~Khot, \emph{On the power of unique 2-prover 1-round games}, STOC 2002, 767--775.
\bibitem{KhotRegev2008}
S.~Khot and O.~Regev, \emph{Vertex cover might be hard to approximate to within $2-\epsilon$}, Journal of Computer and System Sciences 74 (2008), 335--349.
\bibitem{AustrinKhotSafra2011}
P.~Austrin, S.~Khot and M.~Safra, \emph{Inapproximability of vertex cover and independent set in bounded degree graphs}, Theory of Computing 7 (2011), 27--43.
\bibitem{Chan2016}
S.~O. Chan, \emph{Approximation resistance from pairwise-independent subgroups}, Journal of the ACM 63 (2016), Article 27. \url{https://doi.org/10.1145/2873054}.
\bibitem{Halperin2002}
E.~Halperin, \emph{Improved approximation algorithms for the vertex cover problem in graphs and hypergraphs}, SIAM Journal on Computing 31 (2002), 1608--1623.
\bibitem{Trevisan}
L.~Trevisan, \emph{Max Cut and the smallest eigenvalue}, STOC 2009. \url{https://theory.stanford.edu/~trevisan/pubs/maxcutstoc.pdf}.
\bibitem{BermanFujito1999}
P.~Berman and T.~Fujito, \emph{On approximation properties of the independent set problem for low degree graphs}, Theory of Computing Systems 32 (1999), 115--132. \url{https://doi.org/10.1007/s002240000113}.
\bibitem{BrakensiekHuangZwick2024}
J.~Brakensiek, N.~Huang, and U.~Zwick, \emph{Tight approximability of MAX 2-SAT and relatives, under UGC}, Proceedings of SODA 2024, 1328--1344. \url{https://doi.org/10.1137/1.9781611977912.53}.
\bibitem{FeigeKarpinskiLangberg2002}
U.~Feige, M.~Karpinski, and M.~Langberg, \emph{Improved approximation of Max-Cut on graphs of bounded degree}, Journal of Algorithms 43 (2002), 201--219.
\bibitem{HalperinLivnatZwick2004}
E.~Halperin, D.~Livnat, and U.~Zwick, \emph{MAX CUT in cubic graphs}, Journal of Algorithms 53 (2004), 169--185. \url{https://doi.org/10.1016/j.jalgor.2004.06.001}.
\bibitem{HsiehKothari2023}
J.-T. Hsieh and P.~K. Kothari, \emph{Approximating Max-Cut on bounded degree graphs: tighter analysis of the FKL algorithm}, ICALP 2023, arXiv:2206.09204. \url{https://arxiv.org/abs/2206.09204}.
\bibitem{GHLMM2026}
S.~Ghoshal, N.~Huang, E.~Lee, K.~Makarychev, and Y.~Makarychev, \emph{Threshold rounding and bounded-degree Boolean MAX 2-CSP}, arXiv:2607.11050, 2026. \url{https://arxiv.org/abs/2607.11050}.
\bibitem{GoemansWilliamson1995}
M.~X. Goemans and D.~P. Williamson, \emph{Improved approximation algorithms for maximum cut and satisfiability problems using semidefinite programming}, Journal of the ACM 42 (1995), 1115--1145.
\bibitem{KKMO2007}
S.~Khot, G.~Kindler, E.~Mossel, and R.~O'Donnell, \emph{Optimal inapproximability results for MAX-CUT and other 2-variable CSPs?}, SIAM Journal on Computing 37 (2007), 319--357.
\bibitem{MahajanRamesh1999}
S.~Mahajan and H.~Ramesh, \emph{Derandomizing approximation algorithms based on semidefinite programming}, SIAM Journal on Computing 28 (1999), 1641--1663.
\bibitem{OpenAIMath2026}
OpenAI, \emph{Proof of Khot's Unique Games Conjecture} (entry~102 of the OpenAI mathematics release, 6~October 2026), GitHub repository \texttt{openai/math}; not peer reviewed. \url{https://github.com/openai/math}.
\bibitem{RamanRavikumarRao1998}
V.~Raman, B.~Ravikumar, and S.~Srinivasa Rao, \emph{A simplified NP-complete MAXSAT problem}, Information Processing Letters 65 (1998), 1--6.
\bibitem{VegaPalomaRepository}
F.~Vega, \emph{Paloma: MAX-HORN-2SAT and MAX-CUT formalization}, GitHub, 2026. \url{https://github.com/frankvegadelgado/paloma}.
\end{thebibliography}
\end{document}